
\documentclass[a4paper,11pt]{article}
\usepackage[utf8]{inputenc}
\usepackage[T1]{fontenc}
\usepackage[polish]{babel}
\usepackage{listings}
\usepackage{xcolor}
% Konfiguracja wyglądu kodu (opcjonalnie, ale wygląda ładnie)
\lstset{
language=Python,
basicstyle=\ttfamily\small,
keywordstyle=\color{blue},
stringstyle=\color{red},
commentstyle=\color{gray},
showstringspaces=false,
frame=single,
breaklines=true
}
\begin{document}
\section{Wprowadzenie}
W tym dokumencie testowym umieszczamy znane powiedzenie: chon pułk twój i sześć flag.
\section{Kod w Pythonie (metoda lstlisting)}
Poniżej znajduje się skrypt wygenerowany przy użyciu pakietu \texttt{listings}, który jest znacznie potężniejszą meto
\begin{lstlisting}
def powitanie_zaawansowane(nazwa_narzedzia):
print(f"Generuję dokument LaTeX dla: {nazwa_narzedzia}")
status = True
return status
if __name__ == "__main__":
powitanie_zaawansowane("Parser JSON -> LaTeX")
\end{lstlisting}
\end{document}
